\documentclass[a4paper,11pt]{article}

\usepackage{revy}
\usepackage[utf8]{inputenc}
\usepackage[T1]{fontenc}
\usepackage[danish]{babel}


\revyname{MatematikRevyen}
\revyyear{2026}
% HUSK AT OPDATERE VERSIONSNUMMER
\version{0.1}
\eta{$n$ minutter}
\status{Ikke færdig}

\title{Fake-ID}
\author{Forfatter}

\begin{document}
\maketitle

\begin{roles}
\role{G1}[Skuespiller] Rigtig bro-agtig type fra DTU (Seth)
\role{G2}[Skuespiller] Rigtig bro-agtig type fra DTU
\role{G3}[Skuespiller] Alfred, tredje nørd. Snakker ligesom Fogel i \emph{Superbad}
\role{Frivillig}[2 skuespillere] To frivillige, der tjekker studiekort
\role{Dansere}[2 skuespillere] To piger, der danser
\role{Speaker}[Stemme] Speak til sidst
\end{roles}

\begin{props}
\prop{Kæmpe Cafeen-kort}[Person, der skaffer] I stil med det i \emph{Superbad}. Billede, navnet ``Gauss'', fødselsdato i 1643 og en regnbue tegnet med tusch
\prop{DTU-studiekort}[Person, der skaffer] Til G1 og G2
\prop{Bord og stole}[Person, der skaffer] Til de frivillige
\end{props}


\begin{sketch}

\scene{Scenen er delt op i to dele. På venstre side af scenen står to piger og danser. På midten sidder to frivillige og tjekker studiekort. G1 og G2 står uden for.}

\says{G1} Årghh, der er bare så mange damer på Science.

\says{G1} Prøv lige at kig ind ad hegnet.

\scene{G1 og G2 kigger hen på de to, der danser.}

\says{G2} Jaa bro, i mine tre år på Dutten har jeg kun set én pige.

\says{G1} Jeg håber fandme, Alfred har fundet ud af, hvordan han skal komme ind.

\says{G3} Hey, hey, hey bois.

\says{G3} Er I klar??

\says{G2} Fucking Alfred, du gav mig næsten et hjerteanfald.

\says{G2} Har du det fake-ID, vi snakkede om?

\says{G3} Ja ja, to sek, henter det lige.

\scene{G3 hiver et kæmpe Cafeen-kort frem, i stil med det fra \emph{Superbad}.}

\says{G1} Hva fuck er det?

\says{G3} Det' mit Cafeen-kort?

\says{G1} Der er sgu da ikke et billede på et Cafeen-kort!

\says{G3} Alle dem, der laver et fake Cafeen-kort, de kommer jo ikke et billede på.

\says{G3} Derfor valgte jeg at gøre det for at snøre dem, hehe.

\says{G1} Det giver jo ingen mening.

\says{G2} For ikke at snakke om, at du har valgt et dumt fake navn?

\says{G3} Hvad mener du?

\says{G1} Der er sgu da ikke nogen, der hedder Gauss.

\says{G3} Der var da sygt mange, der hed Gauss i vores matematikkurser.

\says{G2} Hvor er du dum!

\says{G2} Det' jo samme person.

\says{G1} For ikke at snakke om, at der står, du er født i 1643, så er du jo over 300 år gammel.

\says{G3} Ja, og der er ingen, der tør at kritisere en person for deres alder.

\says{G3} Desuden er det Newtons fødselsdato.

\says{G1} Hvorfor?!

\says{G3} Fordi hvis de slår Gauss op, finder de Gauss, og så passer det ikke.

\says{G3} Det her kan de ikke slå op.

\says{G2} Og hvorfor har du en regnbue på dit kort?

\says{G3} Det er ikke en regnbue.

\says{G3} Det er et spektrum -- duh.

\says{G3} Et holograms sikkerhedsfeature.

\says{G2} Det er tegnet med tusch.

\says{G3} Ja.

\says{G3} Det tog fire timer.

\says{G1} Kortet er fuldkommen ubrugeligt.

\says{G3} Prøv lige at se her.

\scene{Alfred går hen til de frivillige og kommer igennem uden besvær.}

\says{Frivillig} Næste.

\says{G1} Hvordan gjorde han det?!

\says{G2} Ingen anelse.

\scene{G1 og G2 viser deres DTU-studiekort frem.}

\says{Frivillig} DTU-studiekort virker altså ikke her, dem ser vi hele tiden!

\says{Frivillig} I skal bruge et Cafeen-kort.

\says{G1} Men han havde ikke engang et --

\says{Frivillig} I kan bare købe et her i baren.

\says{Frivillig} 30 kroner.

\says{G1} 30 KRONER?!

\says{G1} Hvad fanden bilder I jer ind?

\says{G1} Hey, det er sgu da ret billigt.

\scene{Segue.}

\says{Speaker} KØB Cafeen-kort!

\says{Speaker} 30 kroner -- gratis kaffe resten af året.

\scene{Tæppe}

\end{sketch}
\end{document}
